\section{Circular Primes --- Problem 35}
\label{prob:35}
\noindent\textbf{Source:} \url{https://github.com/KO0I/projectEuler/blob/reorganize/solutions-by-language/solutions/c/circular_primes_35.c}

\noindent\textbf{Assessment:} Unfinished scaffolding.

\subsection*{Mathematical goal}
A prime is circular when every cyclic rotation of its decimal digits is also prime. Count all circular primes below one million.

\subsection*{How the C code works}
The source describes separate ranges by digit length and sketches functions \verb|a|, \verb|b|, \verb|c| and \verb|d| to examine rotations. The actual \verb|main| sets up a few variables, opens \verb|circprimes|, checks for file-open failure, and closes it. No rotations are calculated and no answer is printed.

\subsection*{Correctness and limitations}
The helper functions have empty bodies despite declaring integer return types. The initialized \verb|sum=9| is neither updated nor printed. A complete solution needs prime membership tests for every rotation, with attention to leading zeros and repeated cyclic forms. The known count is \texttt{...55}, not a result established by this placeholder.


\subsection*{Correction to historical C}
{\color{red}\textbf{Correction:} The empty historical helper stubs have been replaced with a complete sieve and cyclic digit-rotation search for every prime below one million.}
